\section*{Introduction}


  In real world exists several problems which can not be solved by exact methods 
  due to restrictions or constraints, such as time, money, resources, 
  complexity and etc. These problems can be found in almost
  all fields namely engineering, economics, networking, and medical applications
  therefor they need to be handled very carefully to get the best and most efficient solution. 
  As well, it depends on the type of the problem to be solved considering that some of them need to be
  maximized or minimized \cite{talbi2009metaheuristics}. 
  
  For this reason, Metaheuristics Algoritms (MA) were introduced to solve different range of
  problems, particularly non-linear, non-convex, discontinuous, high-dimensional, combinatorial problems
  and whatever problem which can be catagorized as non-deterministic polynomial-time (NP-Hard) 
  \cite{almufti2017using}.

  Every year, numerous MA are developed and inspired by Evolutionary Algorithms (EA), Physical Processes (PP),
  Swarm Intelligence (SI) and Human Behavior (HB) \cite{almufti2022lion}, \cite{agrawal2021metaheuristic}
  However, the most popular and studied MA are Genetic Algorithms (GA) and Simulated Annealing (SA). GA is ispired by
  the adaptation of a living organism in its environment, meanwhile the organisms more adapted would surive
  and reproduce, concurrently the less adapted would extinct, this idea was proposed by Charles Darwin in his book 
  \textit{``The origin of Species''} in 1859 \cite{darwin1859origin}. On the other hand, SA is ispired by... 
  
  % TODO [DO]: Talk about simulated annealing
  % TODO [DO]: Talk about its hybridazation GA + SA   \cite{almufti2022lion} mencions the No free lunch theorem 
  % TODO [DO]: Talk about what my article will contribute

  Therefor this article presents the methodology of GA, SA and GA + SA, a description of the proposed problems,
  the experimental configuration, analysis of the results and finally the conclusions.




  % \Blindtext
